\section*{Bab \ref{ch:b3:electrode-kinetics} --- Kinetika Elektrode dan Elektroanalisis}

\begin{solution}{exo:b3:electrode-kinetics:1}
Kemiringan katodik $2.303RT/\alpha F = \qty{118}{mV}$ memberikan $\alpha = 59.2/118 = 0.50$;
$j_0 = \qty{2.0e-6}{A.cm^{-2}}$, yaitu perpotongan pada $\eta = 0$.
\end{solution}

\begin{solution}{exo:b3:electrode-kinetics:2}
$i_0 = \qty{2.0e-4}{A}$; $R_{\mathrm{ct}} = RT/(Fi_0) = 0.02569/\num{2.0e-4} = \qty{128}{\ohm}$.
\end{solution}

\begin{solution}{exo:b3:electrode-kinetics:3}
Untuk \ce{MgSO4}, $I = \frac12(4c + 4c) = 4c$. Pada \qty{1.0}{mmol/L}, $I = \qty{4.0}{mol.m^{-3}}$
dan $\kappa^{-1} =
\qty{0.96}{nm} \times \sqrt{100/4} = \qty{4.8}{nm}$; pada \qty{0.10}{mol/L}, $I = \qty{400}{mol.m^{-3}}$,
$\kappa^{-1} = \qty{0.48}{nm}$.
\end{solution}

\begin{solution}{exo:b3:electrode-kinetics:4}
$i_{\mathrm c} = C_{\mathrm{dl}}Av = \num{20e-6} \times 0.0707 \times 0.1 = \qty{0.14}{\micro A}$, kurang dari
1\,\% puncaknya. Pada \qty{10}{V/s} nilainya \qty{14}{\micro A}, sedangkan
puncaknya hanya tumbuh menjadi $19\sqrt{100} = \qty{190}{\micro A}$: arus
pengisian tumbuh sebanding dengan $v$, arus faradaik sebanding dengan $\sqrt v$.
\end{solution}

\begin{solution}{exo:b3:electrode-kinetics:5}
$D = \pi\bigl(i\sqrt t/nFAc^*\bigr)^2 = \pi(\num{12.2e-6}/(96\,485 \times 0.0707 \times \num{1.00e-6}))^2 =
\qty{1.0e-5}{cm^2.s^{-1}}$.
\end{solution}

\begin{solution}{exo:b3:electrode-kinetics:6}
$FAc^* = \qty{1.364e-2}{C.cm^{-1}}$, $F/RT = \qty{38.92}{V^{-1}}$. Pada \qty{50}{mV/s}:
$\sqrt{38.92 \times 0.05 \times \num{7.0e-6}} =
\num{3.69e-3}$, $i_{\mathrm p} = 0.4463 \times \num{1.364e-2} \times \num{3.69e-3} = \qty{22.5}{\micro A}$;
pada \qty{500}{mV/s}, $\sqrt{10}$ kali lebih besar, \qty{71}{\micro A}.
\end{solution}

\begin{solution}{exo:b3:electrode-kinetics:7}
(a) Reversibel. (b) \omterm{def:b3:electrode-kinetics:reversibility}{Kuasireversibel}: pemisahan puncaknya tumbuh bersama laju
pindai. (c) Sebuah tahap kimia menyusul transfer elektron (EC): pada pemindaian
lambat, produknya terpakai sebelum sapuan balik; pemindaian yang cepat
mendahului reaksi kimianya.
\end{solution}

\begin{solution}{exo:b3:electrode-kinetics:8}
$K_{ij}a_j/a_i = \num{2e-4} \times 0.14/\num{4.0e-3} = 0.007$: elektrode itu membaca 0.7\,\% terlalu tinggi.
\end{solution}

\begin{solution}{exo:b3:electrode-kinetics:9}
Kuadrat terkecil: $b = \qty{0.1093}{\micro A.L.\micro g^{-1}}$, $a = \qty{0.466}{\micro A}$,
$s = \qty{0.011}{\micro A}$; $c_0 = a/b =
\qty{4.26}{\micro g.L^{-1}}$, dengan $u(c_0) = \frac sb\sqrt{\frac1n + \frac{\bar y^2}{b^2S_{xx}}} = \qty{0.12}{\micro g.L^{-1}}$
($n = 4$, $\bar y = 1.286$, $S_{xx} = 125$).
\end{solution}

\begin{solution}{exo:b3:electrode-kinetics:10}
$Q = 0.0500 \times 1800 = \qty{90.0}{C}$; $n(\ce{Cu}) = Q/2F = \qty{4.66e-4}{mol}$, \qty{29.6}{mg}.
Hasilnya hanya bergantung pada muatan, yang diukur dari arus dan waktu, dan
pada tetapan Faraday: tidak diperlukan standar, asalkan elektrolisisnya
sempurna dan efisiensi arusnya 100\,\%.
\end{solution}

\begin{solution}{exo:b3:electrode-kinetics:11}
$\sqrt{DFv/RT} = \sqrt{\num{1.0e-5} \times 38.92 \times 0.1} = \qty{6.24e-3}{cm.s^{-1}}$:
$\Lambda = 1.6$, \omterm{def:b3:electrode-kinetics:reversibility}{kuasireversibel}. Pada \qty{10}{V/s}, $\Lambda = 0.16$, masih
\omterm{def:b3:electrode-kinetics:reversibility}{kuasireversibel} tetapi lebih jauh dari batas reversibel; gelombangnya baru
akan reversibel di bawah sekitar \qty{1}{mV/s}.
\end{solution}

\begin{solution}{exo:b3:electrode-kinetics:12}
$s/b = \qty{0.0838}{mM}$ dan $b^2S_{xx} = \qty{203.5}{\micro A^2}$. Pada $\bar y$:
$0.0838\sqrt{1 + 1/7} = \qty{0.090}{mM}$. Pada \qty{18.0}{\micro A}:
$0.0838\sqrt{1.143 + 83.0/203.5} = \qty{0.104}{mM}$. Tiga pembacaan pada $\bar y$:
$0.0838\sqrt{1/3 + 1/7} = \qty{0.058}{mM}$.
\end{solution}

\begin{solution}{pb:b3:electrode-kinetics:1}
\textbf{1.} $\ce{C6H12O6 -> C6H10O6 + 2 H+ + 2 e-}$ dan $\ce{[Fe(CN)6]^3- + e- -> [Fe(CN)6]^4-}$;
reaksi keseluruhan
\[
  \ce{C6H12O6 + 2 [Fe(CN)6]^3- -> C6H10O6 + 2 [Fe(CN)6]^4- + 2 H+} .
\]
Di elektrode, $\ce{[Fe(CN)6]^4- -> [Fe(CN)6]^3- + e-}$.
\textbf{2.} Dua.
\textbf{3.} Oksigen terlarut dalam darah sedikit dan berubah-ubah, dan
produknya, hidrogen peroksida, baru teroksidasi pada potensial tinggi tempat
banyak spesies lain ikut bereaksi; mediator tersedia dalam jumlah berlebih
yang diketahui.
\textbf{4.} Agar konsentrasi permukaannya nol: arusnya lalu hanya dibatasi
oleh difusi dan tidak peka terhadap perubahan kecil potensial.
\textbf{5.} Arusnya turun sebanding dengan $t^{-1/2}$; pembacaan harus dilakukan
pada waktu yang dipakai untuk kalibrasi.
\textbf{6.} Cottrell: $i\sqrt t = 42.0$, 41.4, 41.7, 41.6, 41.8: tetap dalam
batas 1\,\%.
\textbf{7.} $D = \pi(\text{kemiringan}/FAc)^2 = \pi(\num{41.8e-6}/(96\,485 \times 0.0300 \times \num{1.00e-5}))^2 = \qty{6.55e-6}{cm^2.s^{-1}}$.
\textbf{8.} $\sqrt{\pi \times \num{6.55e-6} \times 5} = \qty{0.010}{cm} = \qty{100}{\micro m}$:
pengurasan itu mencapai bagian atas lapisan sekitar \qty{5}{s}; pembacaan
yang lebih lambat akan jatuh di bawah hukum Cottrell.
\textbf{9.} Setengahnya: \qty{9.35}{\micro A}.
\textbf{10.} Arus latar: oksidasi pengganggu dan pengotor, pengisian, dan
heksasianoferat(II) sisa milik mediator itu sendiri.
\textbf{11.} Ya: sisaannya, paling besar \qty{0.095}{\micro A}, tersebar tanpa
kecenderungan. Pada glukosa yang tinggi, mediator, yang jumlahnya terbatas,
atau \omterm{def:b3:complex-kinetics:enzyme}{enzimnya} menjadi jenuh dan garisnya melengkung.
\textbf{12.} $(7.10 - 0.356)/0.9189 = \qty{7.34}{mM}$.
\textbf{13.} $0.0838\sqrt{1 + 1/7 + (7.10 - 8.89)^2/203.5} = 0.0838 \times 1.076 = \qty{0.090}{mM}$.
\textbf{14.} Suku $1/m$ (satu pembacaan saja); pembacaan berulang, atau
beberapa strip, mengecilkannya.
\textbf{15.} $3 \times 0.077/0.919 = \qty{0.25}{mM}$.
\textbf{16.} $7.34 \times 180.16 = \qty{1322}{mg.L^{-1}} = \qty{132}{mg.dL^{-1}}$.
\textbf{17.} Askorbat menambahkan arus oksidasinya sendiri: alat ukur membaca
terlalu tinggi.
\textbf{18.} \omterm{def:b3:electrode-kinetics:cv}{Voltamogram} strip tanpa glukosa memperlihatkan gelombang anodik
askorbat di tempat heksasianoferat(II) teroksidasi.
\textbf{19.} Lebih sedikit zat yang teroksidasi pada potensial yang lebih
rendah.
\textbf{20.} Arus blangko dikurangkan dari pembacaan elektrode berenzim
sebelum garis kalibrasi diterapkan.
\textbf{21.} $D$ bertambah beberapa persen per kelvin dan arus sebanding dengan
$\sqrt D$: alat ukur mengukur suhu dan mengoreksinya.
\textbf{22.} Viskositas dan fraksi sel darah merah mengubah $D$ dan arusnya;
standar dalam matriks yang serupa meniadakan pengaruh-pengaruh ini.
\textbf{23.} Variasi antarstrip (luas elektrode, muatan \omterm{def:b3:complex-kinetics:enzyme}{enzim} dan mediator),
suhu, dan komposisi darah menambah ketidakpastian kalibrasi itu sendiri.
\textbf{24.} \textbf{$c(\text{glukosa}) = 7.34 \pm \qty{0.09}{mmol.L^{-1}}$}
(ketidakpastian baku), sekitar \qty{132}{mg.dL^{-1}}.
\end{solution}
